\section*{Bab \ref{ch:g12:synthesis-strategy} --- Strategi Sintesis dan Kimia Hijau}

\begin{solution}{exo:g12:synthesis-strategy:1}
Hidrasi: $46.0/(28.0 + 18.0) = \qty{100}{\percent}$. Fermentasi:
$2 \times 46.0/180.0 = \qty{51.1}{\percent}$.
\end{solution}

\begin{solution}{exo:g12:synthesis-strategy:2}
Kedua gugus, \ce{Z} pada \omterm{def:g11:functional-groups:amine}{amina} glisina dan \omterm{def:g11:functional-groups:carboxylic-acid}{ester} metil pada asam
alanina, dipasang pada tahap 1 dan dilepas pada tahap 3.
\end{solution}

\begin{solution}{exo:g12:synthesis-strategy:3}
Prinsip 5: pakai \omterm{def:g6:solutions-and-solubility:solute}{pelarut} dan kondisi reaksi yang lebih aman.
\end{solution}

\begin{solution}{exo:g12:synthesis-strategy:4}
Etanol berlebih membuat $Q < K$, sehingga reaksinya berlangsung lebih
jauh ke arah maju. \omterm{def:g12:catalysis:catalyst}{Katalis} mempercepat kedua arah secara sama: \omterm{def:g10:reaction-progress-table:system}{keadaan
akhir} yang sama tercapai, hanya lebih cepat.
\end{solution}

\begin{solution}{exo:g12:synthesis-strategy:5}
Ya: hanya \omterm{def:g11:functional-groups:carbonyl}{ketonnya} yang direduksi, menjadi gugus \omterm{def:g11:functional-groups:alcohol}{alkohol} sekunder;
\omterm{def:g11:functional-groups:carboxylic-acid}{esternya} tidak berubah.
\end{solution}

\begin{solution}{exo:g12:synthesis-strategy:6}
\omterm{def:g11:functional-groups:carboxylic-acid}{Ester} lebih banyak: \omterm{def:g7:chemical-reactions:reactant}{pereaksi} berlebih, atau mengambil air (perangkap
air) selagi terbentuk. Lebih cepat: panaskan dengan refluks, tambahkan
\omterm{def:g12:catalysis:catalyst}{katalis} asam.
\end{solution}

\begin{solution}{exo:g12:synthesis-strategy:7}
Dengan \ce{NaBH4}:
\[ \ce{CH3-CH(OH)-CH2-CH2-CO-O-CH3} . \]
Dengan \ce{LiAlH4}:
\[ \ce{CH3-CH(OH)-CH2-CH2-CH2OH} \]
(pentana-1,4-diol), ditambah metanol dari separuh \omterm{def:g11:functional-groups:carboxylic-acid}{ester} yang lain.
\end{solution}

\begin{solution}{exo:g12:synthesis-strategy:8}
\omterm{def:g12:curly-arrows:categories}{Adisi} (\qty{100}{\percent}), esterifikasi (\qty{83}{\percent}), dan
jalur tiga tahap menuju ibuprofen (\qty{77}{\percent}).
$77.4/40.0 = 1.9$: hampir dua kali lebih baik.
\end{solution}

\begin{solution}{exo:g12:synthesis-strategy:9}
$180.0/(138.0 + 102.0) = \qty{75.0}{\percent}$. Asam etanoat:
$60.0/180.0 = \qty{0.333}{kg}$ per kilogram aspirin.
\end{solution}

\begin{solution}{exo:g12:synthesis-strategy:10}
Pemanasan mempercepat reaksi; pendingin mengembalikan uap ke dalam labu,
sehingga tidak ada \omterm{def:g7:chemical-reactions:reactant}{pereaksi} atau \omterm{def:g6:solutions-and-solubility:solute}{pelarut} yang hilang. Pemanasan
meningkatkan laju, bukan \omterm{def:g10:synthesis-yield:yield}{rendemen} akhir.
\end{solution}

\begin{solution}{exo:g12:synthesis-strategy:11}
Enam tahap dibandingkan tiga. Prinsip: mencegah limbah (1),
memaksimalkan \omterm{def:g12:synthesis-strategy:atom-economy}{ekonomi atom} (2), memakai \omterm{def:g12:catalysis:catalyst}{katalis} (9); juga tahap yang
lebih sedikit, energi yang lebih sedikit (6).
\end{solution}

\begin{solution}{exo:g12:synthesis-strategy:12}
$0.80 \times 0.70 \times 0.90 = 0.504$: \qty{50.4}{\percent}. Bahan
awal: $0.10/0.504 = \qty{0.20}{mol}$.
\end{solution}

\begin{solution}{exo:g12:synthesis-strategy:13}
Lindungi \omterm{def:g11:functional-groups:amine}{amina} alanina (\ce{Z}) dan asam glisina (\omterm{def:g11:functional-groups:carboxylic-acid}{ester} metil);
bentuklah ikatan \omterm{def:g11:functional-groups:amine}{amida} antara asam alanina yang bebas dan \omterm{def:g11:functional-groups:amine}{amina} glisina
yang bebas; lepaskan kedua \omterm{def:g12:synthesis-strategy:protecting-group}{gugus pelindung}. Tiga tahapan, yaitu lima
reaksi jika setiap proteksi dan deproteksi dihitung.
\end{solution}

\begin{solution}{exo:g12:synthesis-strategy:14}
\omterm{def:g12:catalysis:catalyst}{Katalis} menerapkan prinsip 9. Tetapi \qty{250}{\celsius} bertentangan
dengan prinsip 6, benzena dengan prinsip 3 dan 5, dan \omterm{def:g12:synthesis-strategy:atom-economy}{ekonomi atom}
\qty{35}{\percent} dengan prinsip 2: $65/35 = \qty{1.9}{kg}$ limbah per
kilogram produk. Satu ciri hijau tidak membuat sebuah proses menjadi
hijau.
\end{solution}

\begin{solution}{exo:g12:synthesis-strategy:15}
% ledger: crit:CO2-T, crit:CO2-p
\qty{40}{\celsius} dan \qty{100}{bar} berada di atas titik kritis
karbon dioksida, \qty{31}{\celsius} dan \qty{73.8}{bar}. Tekanannya
seratus kali tekanan atmosfer: bejananya terbuat dari baja tebal. Karbon
dioksida tidak beracun dan tidak mudah menyala, kembali menjadi gas tanpa
meninggalkan residu, dan dapat dipakai ulang; tidak ada \omterm{def:g6:solutions-and-solubility:solute}{pelarut}
terklorinasi yang lepas.
\end{solution}

\begin{solution}{pb:g12:synthesis-strategy:1}
% ledger: ibu:bhc-ae-recovered
\textbf{1.} C: $10 + 4 + 1 = 15 = 13 + 2$; H: $14 + 6 + 2 = 22 = 18 + 4$;
O: $3 + 1 = 4 = 2 + 2$.

\textbf{2.} Asam etanoat, \ce{C2H4O2}.

\textbf{3.} Nitrogen, klor, dan natrium: ibuprofen hanya mengandung
karbon, hidrogen, dan oksigen.

\textbf{4.} Jalur yang lebih baru: prinsip 9, \omterm{def:g12:catalysis:catalyst}{katalis} alih-alih \omterm{def:g7:identifying-substances:characteristic-test}{reagen}
stoikiometris.

\textbf{5.} $13 \times 12.0 + 18 \times 1.0 + 2 \times 16.0 =
\qty{206.0}{g/mol}$.

\textbf{6.} $134.0 + 102.0 + 122.5 + 68.0 + 19.0 + 33.0 + 36.0 =
\qty{514.5}{g/mol}$.

\textbf{7.} $206.0/514.5 = \qty{40.0}{\percent}$.

\textbf{8.} $134.0 + 102.0 + 2.0 + 28.0 = \qty{266.0}{g/mol}$.

\textbf{9.} $206.0/266.0 = \qty{77.4}{\percent}$.

\textbf{10.} $(206.0 + 60.0)/266.0 = \qty{100}{\percent}$ di atas kertas;
lebih dari \qty{99}{\percent} dalam praktik.

\textbf{11.} $514.5/206.0 = \qty{2.50}{t}$.

\textbf{12.} $2.50 - 1 = \qty{1.50}{t}$.

\textbf{13.} $266.0/206.0 = \qty{1.29}{t}$.

\textbf{14.} $60.0/206.0 = \qty{0.29}{t}$.

\textbf{15.} $1.50 - 0.29 = \qty{1.21}{t}$.

\textbf{16.} $0.90^6 = 0.53$: \qty{53}{\percent}.

\textbf{17.} $0.90^3 = 0.73$: \qty{73}{\percent}.

\textbf{18.} Setiap tahap memakai \omterm{def:g6:solutions-and-solubility:solute}{pelarut}, pemisahan, dan pemurnian,
serta kehilangan sebagian produk: tahap yang lebih sedikit berarti
kehilangan semacam itu juga lebih sedikit.

\textbf{19.} Hampir tidak ada: satu-satunya hasil sampingnya dipakai.

\textbf{20.} Jalur tiga tahap menghindari sekitar \qty{1.5}{t} limbah per
ton ibuprofen (\qty{1.2}{t} bahkan jika asam etanoatnya dibuang).
\end{solution}
